% answers_short.tex —— 1956 年·全国卷 · 简短答案速查
% 由 tools/make_short_answers.py 自 answers.tex 的 \gksolhead/\gkanswer 提取，
% 并与 1956/1956_数学真题_答案要点.md 逐题比对后生成，勿手改（改源 MD 后重跑）。

\begin{gknotice}
\textbf{使用说明}：本节只给最终答案便于快速核对，完整推导见解析版。答案为编写者独立推导，\textbf{非官方标准答案}。
\end{gknotice}

\gkpart{小题}
\begin{itemize}[leftmargin=2.6em]
  \item[\textbf{1.}] $1$
  \item[\textbf{2.}] 对任意实数 $m$，方程有两个不相等的实根。
  \item[\textbf{3.}] $S_{\triangle AEF}=\dfrac12S_{\triangle ABC}$。
  \item[\textbf{4.}] $120^\circ$。
  \item[\textbf{5.}] $\sin(\alpha+\beta)=\cos(\alpha+\beta)$。
\end{itemize}

\gkpart{大题}
\begin{itemize}[leftmargin=2.6em]
  \item[\textbf{6.}] $(x,y)=\left(\frac{9+\sqrt{191}}2,\frac{9-\sqrt{191}}2\right),\ \left(\frac{9-\sqrt{191}}2,\frac{9+\sqrt{191}}2\right),\ (10,6),\ (6,10)$
  \item[\textbf{7.}] $PA^2=AB^2+PB\cdot PC$
  \item[\textbf{8.}] 平面 $A'ACC'$ 垂直于底面 $ABCD$
  \item[\textbf{9.}] 三个角中间的一个为 $60^\circ$，所以必有一个角是 $60^\circ$
  \item[\textbf{10.}] 三边相等，故三个角都是 $60^\circ$
\end{itemize}
